% Gerado por scripts/etapa5_transbordo.py — requer \usepackage{booktabs}
\begin{table}[htbp]
\centering
\small
\caption{Tamanho da aproximação do alvo de t+2 pelo de t+1}
\label{tab:tb_aproximacao_alvo}
\begin{tabular}{lrrrr}
\toprule
Nível & Alvo médio & Diferença com sinal & Diferença absoluta & Pares com alvo diferente \\
\midrule
80\% & 2,960 & 0,3\% & 10,9\% & 23,0\% \\
90\% & 3,981 & 0,2\% & 8,1\% & 23,0\% \\
95\% & 4,995 & 0,2\% & 6,5\% & 23,1\% \\
98\% & 6,319 & 0,1\% & 5,3\% & 23,8\% \\
99\% & 7,270 & 0,1\% & 4,7\% & 24,0\% \\
\bottomrule
\end{tabular}
\par\smallskip
{\footnotesize Nota: simulação de 90 dias (26/01 a 24/04/2016), combinação, 1.437 itens × 4 lojas; reposição diária pelo fornecedor até o quantil do nível de serviço, arredondado para o inteiro mais próximo; venda perdida. Diferença entre o alvo de t+2 e o de t+1 (o usado como aproximação no prazo de 2 dias), somada sobre séries e dias e dividida pela soma dos alvos de t+1. Com sinal: positivo = alvo de t+2 maior (a aproximação subestima a necessidade). Pares com alvo diferente: \% dos pares de dias consecutivos com alvos distintos.}
\end{table}
